% Compile twice with pdflatex. All references are included below.
\documentclass[11pt,a4paper]{article}
\usepackage[T1]{fontenc}
\usepackage{lmodern,amsmath,amssymb,amsthm,microtype}
\usepackage[margin=28mm,top=25mm,bottom=25mm]{geometry}
\usepackage{xcolor}
\definecolor{linkink}{HTML}{164C66}
\usepackage[colorlinks=true,linkcolor=linkink,citecolor=linkink,urlcolor=linkink,
  pdftitle={A minimal action of two commuting maps with no semi-open nonidentity element},
  pdfsubject={Topological dynamics of commutative semigroup actions}]{hyperref}
\usepackage{fancyhdr}
\pagestyle{fancy}\fancyhf{}\renewcommand{\headrulewidth}{0pt}
\fancyfoot[L]{\footnotesize\textsf{Semi-openness in minimal semigroup actions}}
\fancyfoot[R]{\footnotesize\thepage}
\setlength{\footskip}{13mm}
\setlength{\parindent}{1.15em}
\newcommand{\T}{\mathbb{T}}
\newcommand{\N}{\mathbb{N}}
\newcommand{\NN}{\mathbb{N}_0}
\newcommand{\inter}{\operatorname{Int}}
\newtheorem{theorem}{Theorem}
\newtheorem{lemma}[theorem]{Lemma}
\newtheorem{corollary}[theorem]{Corollary}
% Add an author line here only if authorship is established.
\begin{document}
\begin{center}
{\LARGE\bfseries A minimal action of two commuting maps\\[2pt]
with no semi-open nonidentity element\par}
\vspace{6pt}
{\small 20 November 2026}
\end{center}

\begin{abstract}\noindent
We construct a minimal action of $\NN^2$ by continuous surjections on a compact
metrizable continuum. For each nonzero $(m,n)$, the corresponding map sends
a nonempty open set onto the same closed nowhere dense slice. The action is a
factor of a minimal action whose maps are all open. The construction answers
the commutative-semigroup question of Blahodatna, Detmold, Glasscock and Le
and gives a negative answer to the metrizable case of a question of Dai,
Feng, Lv and Xie on factor maps of minimal semiflows.
\end{abstract}

\section{The question}
A continuous map $f:X\to Y$ is \emph{semi-open} if $\inter_Y f(U)\ne\varnothing$
for every nonempty open $U\subseteq X$. An action of a discrete semigroup $S$
is \emph{minimal} if $\overline{Sx}=X$ for every $x\in X$.
A factor map is a continuous surjection intertwining the two actions.

Kolyada, Snoha and Trofimchuk proved that a minimal continuous self-map of
a compact Hausdorff space is semi-open \cite[Theorem 2.4(1)]{KST}. Its
iterates are then semi-open by composition. For actions with several
generators, Blahodatna, Detmold, Glasscock and Le gave a noncommutative
example and asked whether a commutative example exists
\cite[\S1.2, footnote~1]{BDGL}.

\begin{theorem}\label{thm:main}
There exist a compact connected metrizable space $X$, commuting continuous
surjections $s,r:X\to X$, and a closed nowhere dense set $D\subset X$ such
that the $\NN^2$-action $(m,n)\cdot x=s^m r^n x$ is minimal.
For every $m+n>0$ there exists a nonempty open set $U_{m,n}\subset X$
with $s^m r^n(U_{m,n})=D$.
\end{theorem}

Akin, Auslander and Nagar used a one-sided shift and a coordinate map
constant on an interval to produce a non-semi-open map
\cite[Example 2.30]{AAN}. In their construction the shift and the
coordinatewise map preserve the set of constant sequences. To obtain
joint minimality here, we first construct a skew translation with a dense
orbit in one entire fiber. We then pass to a factor that retains only
partial information about the first sequence coordinate.

\section{A minimal action by open maps}
Write $\T=\mathbb R/\mathbb Z$, $\NN=\{0,1,\ldots\}$ and
$\N=\{1,2,\ldots\}$. All products carry the product topology.
Set $c(t)=\cos(2\pi t)$. Choose $\alpha\in\T$ irrational so that
\begin{equation}\label{eq:ind}
1,\ c(\alpha),\ c(2\alpha),\ c(3\alpha),\ldots
\quad\text{are linearly independent over }\mathbb Q.
\end{equation}
Such a choice exists. Each nonzero finite integer relation prescribes the
zeros of a nonzero trigonometric polynomial
$a_0+\sum_{j=1}^{k}a_j\cos(2\pi jt)$. Its Fourier coefficients show that
it is nonzero, and multiplication by $z^k$, $z=e^{2\pi i t}$, turns it into
an ordinary polynomial. Its zero set on $\T$ is finite. Excluding these
countably many finite sets and the rational points gives \eqref{eq:ind}.

Put $K=\T^{\NN}$ and $Z=\T\times K$. For $x=(x_j)_{j\ge0}$ define
\begin{align*}
 A(t,x)&=\bigl(t+\alpha,(x_{j+1})_{j\ge0}\bigr),\\
 H(t,x)&=\bigl(t,(x_j+c(t+j\alpha))_{j\ge0}\bigr), \qquad B=H A .
\end{align*}
The additions in sequence coordinates are taken modulo one. In the
$j$th coordinate both $AH$ and $HA$ give $x_{j+1}+c(t+(j+1)\alpha)$.
Thus $A$ and $B$ commute. The unilateral shift sends basic cylinders to
open cylinders and is onto. Hence $A$ is open and onto, and so is $B$,
since $H$ is a homeomorphism with inverse given by subtracting the same
coordinate increments.

\begin{lemma}\label{lem:min}
The forward action generated by $A$ and $B$ on $Z$ is minimal.
\end{lemma}
\begin{proof}
Every nonempty closed forward-invariant set contains a minimal such set
$M$, by compactness and the nested-intersection argument. By commutativity,
$A(M)$ and $B(M)$ are nonempty compact forward-invariant subsets of $M$.
Consequently $A(M)=M=B(M)$ and
\[
 H(M)=H(A(M))=B(M)=M.
\]
The first-coordinate projection of $M$ is closed and invariant under the
irrational rotation $t\mapsto t+\alpha$, so it is all of $\T$. In particular,
$M$ meets the fiber over $\alpha$. On this fiber $H$ translates by the
vector $(c((j+1)\alpha))_{j\ge0}$. By \eqref{eq:ind} and Kronecker's
theorem \cite[Theorem 2.39]{EFHN}, every forward orbit of this translation
is dense in every finite initial coordinate block. These blocks determine
the product topology. Thus the orbit of any point of $M$ on this fiber
is dense in $\{\alpha\}\times K$, and $M$ contains the whole fiber.

As the shift is onto, $A$ takes each whole fiber onto the next one.
It follows that $M$ contains $\{k\alpha\}\times K$ for every $k\ge1$.
Their union is dense in $Z$, so $M=Z$. Applying this conclusion to each
forward orbit closure proves the lemma.
\end{proof}

\section{A factor with a nowhere dense image}
Let $p:\T\to[0,1]$ be the continuous surjection $p(v)=\max\{0,c(v)\}$.
It vanishes on the open arc $J=(1/4,3/4)\subset\T$. Define
\[
 X=\T\times[0,1]\times\T^{\N},\qquad
 q(t,x_0,x_1,\ldots)=(t,p(x_0),x_1,x_2,\ldots).
\]
The countable product $X$ is compact, metrizable and path-connected.
Write its points as $(t,u,x_1,x_2,\ldots)$. Both $A$ and $B$ discard
$x_0$. It follows that $qA$ and $qB$ are constant on each fiber of $q$.
Since $q$ is a continuous surjection from a compact space to a Hausdorff
space, it is a quotient map. There are therefore continuous maps $s,r$ on
$X$ with $sq=qA$ and $rq=qB$. They are surjective and commute because
$A,B$ are surjective and commute and $q$ is onto. Their joint action is
minimal by Lemma~\ref{lem:min}, since the image of each dense orbit is dense.

\begin{proof}[Proof of Theorem~\ref{thm:main}]
The first generator has the formula
\[
 s(t,u,x_1,x_2,\ldots)=(t+\alpha,p(x_1),x_2,x_3,\ldots).
\]
In particular it sends the cylinder $x_1\in J$ onto
$D=\T\times\{0\}\times\T^{\N}$. The same slice works for every nonzero
element. Fix $d=m+n>0$. Upstairs, $A^mB^n=H^n A^d$, so downstairs
the value of $s^m r^n$ has base coordinate $t+d\alpha$, interval coordinate
$p(x_d+n c(t+d\alpha))$, and tail coordinates
\[
 x_{d+j}+n c(t+(d+j)\alpha),\qquad j\ge1.
\]
Therefore the open set
\[
 U_{m,n}=\{(t,u,x_1,\ldots)\in X:\ x_d+n c(t+d\alpha)\in J\}
\]
maps into $D$. To reach any prescribed point $(t',0,y_1,y_2,\ldots)$ of
$D$, set $t=t'-d\alpha$, choose $x_d$ so that $x_d+n c(t')\in J$, and set
$x_{d+j}=y_j-n c(t'+j\alpha)$ for $j\ge1$. All other input coordinates
are free. This proves $s^m r^n(U_{m,n})=D$ and also proves nonemptiness of
$U_{m,n}$. The set $D$ is closed with empty interior because $\{0\}$
has empty interior in $[0,1]$.
\end{proof}

\section{Two consequences}
\begin{corollary}
A continuous equivariant surjection between compact metrizable minimal
$\NN^2$-systems can fail to be semi-open even when it is a self-map.
\end{corollary}
\begin{proof}
Take $\pi=s:X\to X$. The map is onto, and it intertwines every element of
the action since $s$ and $r$ commute. Theorem~\ref{thm:main} with $(m,n)=(1,0)$
gives the required open set.
\end{proof}
This answers the metrizable branch of \cite[Question 5.2.4]{DFLX}. The
fiber-onto hypothesis of their positive theorem \cite[Theorem 3.2]{DFLX}
fails here in an explicit way. It would require
$g(\pi^{-1}(y))=\pi^{-1}(g y)$ for all $g$ and $y$. Taking $g=\pi=s$,
the left side is $\{y\}$, whereas the right side contains every point
obtained from $y$ by varying its interval coordinate.

\begin{corollary}
The property that all action maps are open can be lost on taking a factor
of a compact metrizable minimal commutative action. In that factor all
nonidentity action maps can even fail to be semi-open.
\end{corollary}
\begin{proof}
Every $A^mB^n$ on $Z$ is a composition of open maps. The quotient $q$
intertwines these maps with those of Theorem~\ref{thm:main}.
\end{proof}

\paragraph{Provenance.}
OpenAI's ChatGPT generated the counterexample and an initial proof in
response to the research question and assisted in preparing this manuscript.

\begin{thebibliography}{5}\small
\bibitem{AAN} E. Akin, J. Auslander and A. Nagar,
\emph{Variations on the concept of topological transitivity},
Studia Math. \textbf{235} (2016), 225--249.
\href{https://doi.org/10.4064/sm8553-7-2016}{doi:10.4064/sm8553-7-2016}.
\bibitem{BDGL} A. Blahodatna, L. Detmold, D. Glasscock and A.\,N. Le,
\emph{The regionally proximal relation for commutative semigroup actions},
\href{https://arxiv.org/abs/2610.10397v1}{arXiv:2610.10397v1} (2026).
\bibitem{DFLX} X. Dai, L. Feng, C. Lv and Y. Xie,
\emph{On semi-openness of fiber-onto extensions of minimal semiflows and
quasi-separable maps}, \href{https://arxiv.org/abs/2601.03380v2}{arXiv:2601.03380v2} (2026).
\bibitem{EFHN} T. Eisner, B. Farkas, M. Haase and R. Nagel,
\emph{Operator Theoretic Aspects of Ergodic Theory}, Graduate Texts in
Mathematics 272, Springer, 2015.
\href{https://doi.org/10.1007/978-3-319-16898-2}{doi:10.1007/978-3-319-16898-2}.
\bibitem{KST} S. Kolyada, \v{L}. Snoha and S. Trofimchuk,
\emph{Noninvertible minimal maps}, Fund. Math. \textbf{168} (2001), 141--163.
\href{https://doi.org/10.4064/fm168-2-5}{doi:10.4064/fm168-2-5}.
\end{thebibliography}
\end{document}
